\documentclass[12pt]{amsart}

\usepackage{amsmath,amssymb,amsthm}
\usepackage{hyperref}
\usepackage{enumitem}

%% Theorem environments
\newtheorem{theorem}{Theorem}[section]
\newtheorem{lemma}[theorem]{Lemma}
\newtheorem{corollary}[theorem]{Corollary}
\newtheorem{proposition}[theorem]{Proposition}
\newtheorem{definition}[theorem]{Definition}
\newtheorem{remark}[theorem]{Remark}
\newtheorem{example}[theorem]{Example}

%% Notation macros
\newcommand{\Hsym}{\mathfrak{H}_{\mathrm{sym}}}
\newcommand{\Xsym}{\mathfrak{X}_{\mathrm{sym}}}
\newcommand{\Li}{\mathrm{Li}}
\newcommand{\Ofinite}{O_{\mathrm{fin}}}   % use \Ofinite^{(K)} at call sites
\newcommand{\Otheta}{O_{\theta}^{\mathrm{samp}}}
\newcommand{\Ooracle}{O_{\mathrm{oracle}}}
\newcommand{\re}{\operatorname{Re}}
\newcommand{\im}{\operatorname{Im}}
\newcommand{\NZ}{N_{\mathcal{Z}}}
\newcommand{\CC}{\mathbb{C}}
\newcommand{\RR}{\mathbb{R}}
\newcommand{\QQ}{\mathbb{Q}}
\newcommand{\ZZ}{\mathbb{Z}}
\newcommand{\TT}{\mathbb{T}}
\newcommand{\incomp}{\mathbin{\bowtie}}

\begin{document}

\title[Finite Li Statistics and Critical-Line Placement]{Arithmetic
  Information Barriers for the Riemann Hypothesis}

\author{Lin Tao}
\date{2026}

\begin{abstract}
We prove three theorems about the standard Li-type observation of finite
symmetric zero multisets.
Let $\Hsym$ be the class of finite multisets in the open critical strip
symmetric under $\rho\mapsto\bar\rho$ and $\rho\mapsto 1-\rho$, and let
$P(\mathcal{Z})=1$ iff every atom has $\re\rho=1/2$.

\textbf{Theorem~1} (exact Li collision): for every $m\geq 1$ there exist
$\mathcal{Z}_+,\mathcal{Z}_-\in\Hsym$ with $P(\mathcal{Z}_+)=1$,
$P(\mathcal{Z}_-)=0$, and identical standard Li observations
$\Li_j(\mathcal{Z}_+)=\Li_j(\mathcal{Z}_-)$ for $j=1,\ldots,m$.
The collision is realised by explicit rational-arithmetic witnesses.

\textbf{Theorem~2} (no uniform separation): no fixed finite collection of
standard Li values enforces a strictly positive gap between $P=1$ and $P=0$
within $\Hsym$.

\textbf{Theorem~3} (incomparability): the standard Li observation map
$\Ofinite^{(K)}$ and the sampled counting map $\Otheta$ are incomparable in
the information refinement preorder on $\Hsym$: each can distinguish
inputs the other cannot.

We also record, as a corollary of Andersson~\cite{Andersson2024}, that no
finite Euler-factor constraint on a Helson zeta function forces its zeros
to lie on the critical line.

None of these results assume or imply the Riemann Hypothesis.
The results establish that, within $\Hsym$, neither $\Ofinite^{(K)}$ nor
$\Otheta$ can determine the predicate $P$, and no function $f$ satisfies
$P = f \circ \Ofinite^{(K)}$ on all of $\Hsym$.
\end{abstract}

\maketitle
\tableofcontents

%%--------------------------------------------------------------------
\section{Introduction}
\label{sec:intro}
%%--------------------------------------------------------------------

Li's criterion~\cite{Li1997} asserts that RH is equivalent to the
nonnegativity of all Li coefficients
$\lambda_n = \sum_\rho [1-(1-1/\rho)^n]$,
where the sum is taken in the symmetric sense
$\lim_{T\to\infty}\sum_{|\im\rho|\leq T}$
(as in Bombieri--Lagarias~\cite{BombieriLagarias1999}).
This reformulation raises a natural question: can one establish RH by
verifying, via a uniform template over all $n$, that $\lambda_n \geq 0$?
The results of this paper say that a strategy based on any \emph{fixed
finite} collection of Li values cannot decide the target predicate even on
the simplified class $\Hsym$ of finite symmetric multisets.  A proof scheme
that separately handles each $n$ but uses a uniform method reduces to the
full Li criterion and is not blocked.

\subsection*{What is an information obstruction?}

We work with an \emph{observation map} $O : \Hsym \to S$, where $\Hsym$
is the ambient class of finite symmetric zero multisets and $S$ is a set
of observable records.  An \emph{information obstruction for $O$} is an
explicit pair $(\mathcal{Z}_+, \mathcal{Z}_-) \in \Hsym^2$ such that:
\begin{enumerate}[label=(\roman*)]
  \item $O(\mathcal{Z}_+) = O(\mathcal{Z}_-)$ (exact collision --- $O$
    cannot distinguish the two inputs);
  \item $P(\mathcal{Z}_+) \neq P(\mathcal{Z}_-)$ (the two inputs give
    opposite answers to the target predicate);
  \item neither $\mathcal{Z}_+$ nor $\mathcal{Z}_-$ involves RH as a
    hypothesis.
\end{enumerate}
An obstruction pair shows that no function $f : S \to \{0,1\}$ satisfies
$P(\mathcal{Z}) = f(O(\mathcal{Z}))$ for all $\mathcal{Z}\in\Hsym$: the
observation $O$ alone cannot decide $P$.

\begin{remark}[What these results do and do not say]
\label{rem:scope}
All obstructions are proved for the class $\Hsym$ of \emph{finite}
symmetric multisets, not for the actual zero set of $\zeta$.  The actual
zero set is infinite, satisfies a zero-counting asymptotic, and is
governed by the Hadamard product, functional equation, Gamma factor, and
Euler product --- none of which are imposed on $\Hsym$.  Accordingly, the
theorems say: a method that reads only $O(\mathcal{Z})$ for some $O \in
\{\Ofinite^{(K)}, \Otheta\}$ cannot determine $P$ on $\Hsym$.  They do
not directly bound what a method using the full analytic structure of
$\zeta$ can achieve.

An information obstruction does not prevent RH from being true, nor does it
prevent an RH proof that uses structure not captured by $O$.  The escape
route for each theorem is named explicitly.
\end{remark}

\begin{remark}[Analogy with complexity-theoretic barriers]
\label{rem:RR}
The closest formal analogues in complexity theory are the
Razborov--Rudich natural proofs barrier~\cite{RazborovRudich1997} and
the relativization barrier~\cite{BakerGillSolovay1975}.  The present
framework is \emph{inspired by} the Razborov--Rudich structure in that
it exhibits explicit adversarial witnesses showing a class of observation
maps cannot decide the target predicate.  However, the natural proofs
framework additionally requires constructivity and density hypotheses and
involves a diagonalization against pseudorandom functions; those features
are absent here.  The analogy is structural and motivational, not a
formal reduction.
\end{remark}

\subsection*{Results}

The ambient class $\Hsym$ consists of finite zero multisets in the open
critical strip $\{0 < \re \rho < 1\}$, closed under $\rho \mapsto \bar\rho$
and $\rho \mapsto 1 - \rho$ (Section~\ref{sec:setup}).  The target
predicate is $P(\mathcal{Z}) = 1$ iff every atom has $\re \rho = 1/2$.
The \emph{standard Li tests} are $\varphi_j(\rho) = 1-(1-\rho^{-1})^j$.

\begin{itemize}
\item \textbf{Theorem~1} (Section~\ref{sec:B2}): For every fixed $m\geq 1$,
  there exist $\mathcal{Z}_+, \mathcal{Z}_- \in \Hsym$ with
  $P(\mathcal{Z}_+) = 1$, $P(\mathcal{Z}_-) = 0$, and
  $\Li_j(\mathcal{Z}_+) = \Li_j(\mathcal{Z}_-)$ exactly for
  $j = 1, \ldots, m$.  The collision is realised by explicit
  rational-arithmetic witnesses.

\item \textbf{Theorem~2} (Section~\ref{sec:B1}): No fixed finite collection
  of standard Li values admits a uniform strictly positive lower bound
  separating $P(\mathcal{Z}) = 1$ from $P(\mathcal{Z}) = 0$ within $\Hsym$.

\item \textbf{Theorem~3} (Section~\ref{sec:H}): Fix $M \geq 1$ and
  positive sampling levels $0 < d_1 < \cdots < d_M$.  For every $K \geq 1$,
  the Li observation map $\Ofinite^{(K)}$ and the sampled counting map
  $\Otheta$ are incomparable in the information refinement preorder on
  $\Hsym$.

\item \textbf{Corollary~4} (Section~\ref{sec:C}): For every finite prime
  cutoff $P_0$ and every $z_1 \in S$ with $\re z_1 \neq 1/2$, there
  exists a Helson zeta function that agrees with $\zeta$ on all Euler
  factors $p \leq P_0$ yet has a zero at $z_1$.
\end{itemize}

All four results are unconditional; none assumes RH.

\subsection*{Comparison with prior work}

\textbf{Li's criterion and reformulations.}
Li's criterion~\cite{Li1997} asserts RH is equivalent to $\lambda_n \geq 0$
for all $n$; Bombieri--Lagarias~\cite{BombieriLagarias1999} gave an
equivalent formulation in terms of explicit-formula positivity.  These are
RH-\emph{equivalences}, not obstructions: they locate difficulty but do not
block proof methods.  The distinction is categorical: an equivalence
$C \Leftrightarrow \mathrm{RH}$ says the problem is hard; our results say
that a specific \emph{observation class} cannot decide the target predicate
within $\Hsym$.

\textbf{Weil explicit formula and positivity.}
The Weil explicit formula expresses sums over zeros as sums over primes
plus archimedean contributions; positivity of these sums for all test
functions would imply RH.  Connes~\cite{ConnesWeilFF1999} interpreted this
as a spectral absorption problem on an adelic space, and
Bombieri~\cite{Bombieri2000} reformulated it as a positivity problem on
$L^2$ spaces over the primes.  These approaches use \emph{all} test
functions simultaneously; our results apply precisely when only a
\emph{finite} family is used.

\textbf{Selberg trace formula and spectral methods.}
The Selberg trace formula connects eigenvalues of the Laplacian on
hyperbolic surfaces to geodesic lengths, and analogies with the Riemann
zeta function have been studied extensively (see Hejhal~\cite{Hejhal1976}).
Connes--Consani--Moscovici~\cite{CCM2025} construct zeta spectral triples
whose Dirac-operator regularized determinant encodes $\hat{\xi}(z)$;
establishing the convergence of suitably normalized determinants to the
Riemann $\Xi$ function remains an open step.

\textbf{Helson class and Dirichlet series.}
Helson~\cite{Helson1969} introduced random Dirichlet series with unimodular
coefficients; their zero distribution was studied by
Seip~\cite{Seip2020} and Bochkov--Romanov~\cite{BochkovRomanov2022}.
Andersson~\cite{Andersson2024} proved that prescribed zeros in the critical
strip can be realized by Helson zeta functions; Corollary~4 of this paper
makes the finite-Euler-factor limitation explicit as an information
obstruction.

\textbf{Rational arithmetic witnesses.}
The exact rational witnesses for the $m=2$ case have been verified by
independent computation using exact rational arithmetic
(\texttt{fractions.Fraction} in Python).

%%--------------------------------------------------------------------
\section{Setup and notation}
\label{sec:setup}
%%--------------------------------------------------------------------

\subsection{Ambient class}

\begin{definition}[Ambient class $\Hsym$]
\label{def:ambient}
A finite multiset $\mathcal{Z}$ of complex numbers is in $\Hsym$ if:
\begin{enumerate}[label=\textup{(S\arabic*)}]
  \item\label{itm:strip} every atom satisfies $0 < \re\rho < 1$;
  \item\label{itm:conj} $\rho \in \mathcal{Z}$ implies $\bar\rho \in \mathcal{Z}$
    with the same multiplicity;
  \item\label{itm:fe} $\rho \in \mathcal{Z}$ implies $1-\rho \in \mathcal{Z}$
    with the same multiplicity.
\end{enumerate}
The target predicate is $P(\mathcal{Z}) = 1$ iff $\re\rho = 1/2$ for every
$\rho \in \mathcal{Z}$.
\end{definition}

Condition~\ref{itm:strip} excludes $0$ and $1$, ensuring that
$\varphi_j(\rho)$ is well-defined for all standard Li tests.
Conditions~\ref{itm:conj} and~\ref{itm:fe} mimic the two symmetries of
the nontrivial zeros of~$\zeta$.

\begin{remark}[Finiteness]
We work with \emph{finite} multisets throughout.  This is the natural
domain for the finite-observation theorems~1--3.  Corollary~4 involves an
infinite Helson zeta function, whose zero set in the strip is treated
separately.
\end{remark}

\subsection{Standard Li tests and the finite observation map}

\begin{definition}[Standard Li tests]
\label{def:Listd}
For $j \geq 1$, the \emph{$j$-th standard Li test} is
\[
  \varphi_j(\rho) = 1 - \bigl(1 - \rho^{-1}\bigr)^j.
\]
These satisfy $\varphi_j(\bar\rho) = \overline{\varphi_j(\rho)}$
(reality condition).  For $\sigma\in(0,1)$ and $|t|\geq 1$,
\[
  |\varphi_j(\sigma+it)| = \frac{j}{|\sigma+it|} + O\bigl(|t|^{-2}\bigr)
  = O\bigl(|t|^{-1}\bigr).
\]
In particular, $\varphi_j = O(|t|^{-1})$ uniformly in $\sigma\in(0,1)$,
but not $O(|t|^{-1-\delta})$ for any $\delta>0$.
\end{definition}

\begin{remark}[Other test families]
\label{rem:othertests}
Moment-type tests $\varphi_j(\rho)=\rho^{-k_j}$ admit an analogous
cosine-Vandermonde rank argument (Remark~\ref{rem:momentVander}); the
proof of Theorem~1 uses the Chebyshev--Vandermonde structure specific to
the standard Li tests and does not automatically extend to other families
unless one can verify rationality of the resulting system over $\QQ$.
Weil-type tests $\varphi_h(\rho)=\hat{h}((\rho-\tfrac{1}{2})/i)$ for
$h\in C_c^\infty(\RR)$ are not meromorphic functions of $\rho$ in general;
their treatment is deferred.
\end{remark}

\begin{definition}[Li observation and finite observation map]
\label{def:Li}
For $\mathcal{Z} \in \Hsym$ and $j \geq 1$, define
\[
  \Li_j(\mathcal{Z}) = 2 \sum_{\rho \in \mathcal{Z}} \varphi_j(\rho).
\]
This sum is real: by the reflection symmetry~\ref{itm:fe}, each atom
$\rho$ is paired with $1-\rho$, and $\varphi_j(\rho)+\varphi_j(1-\rho)$
is real by conjugation symmetry~\ref{itm:conj} and the reality condition.

The \emph{finite observation map} with $K$ tests is
\[
  \Ofinite^{(K)} : \Hsym \longrightarrow \RR^K, \qquad
  \Ofinite^{(K)}(\mathcal{Z}) = \bigl(\Li_1(\mathcal{Z}), \ldots, \Li_K(\mathcal{Z})\bigr).
\]
\end{definition}

\begin{remark}[Normalization convention]
\label{rem:normconv}
The factor of $2$ in Definition~\ref{def:Li} is a non-standard convention
(``R-symm'') adopted throughout this paper.  The standard convention
(``R-atom'') sums $\varphi_j$ once per atom, giving a value half as large.
Every qualitative conclusion (existence of collisions, absence of uniform
gaps) is convention-independent; all numerical anchors use R-symm.  The
classical Li coefficients $\lambda_n$ of $\zeta$ are defined via a
symmetric limit $\lim_{T\to\infty}\sum_{|\im\rho|\leq T}$ over the
actual infinite zero set~\cite{BombieriLagarias1999} and should not be
confused with $\Li_j(\mathcal{Z})$ for finite $\mathcal{Z}\in\Hsym$.
\end{remark}

\subsection{Sampled counting map}

\begin{definition}[Sampled counting map]
\label{def:Otheta}
Fix $M \geq 1$ and positive sampling levels $0 < d_1 < \cdots < d_M$.
For $\mathcal{Z} \in \Hsym$, let
$\NZ(u) = \#\{\rho \in \mathcal{Z} : |\im\rho| \leq u\}$
(with multiplicity; this counts both $\rho$ and $\bar\rho$).
The \emph{sampled counting map} is
\[
  \Otheta : \Hsym \longrightarrow \ZZ_{\geq 0}^M, \qquad
  \Otheta(\mathcal{Z}) = \bigl(\NZ(d_1), \ldots, \NZ(d_M)\bigr).
\]
\end{definition}

The map $\Otheta$ is nonconstant: two symmetric quartets with imaginary
parts on opposite sides of some $d_m$ give different values.  It is
$\mathcal{Z}$-dependent, in contrast to the Gram-level sequence
$\theta_{\mathrm{level}}(n)$ used in the spectral companion paper, which
is a fixed, $\mathcal{Z}$-independent sequence.

\subsection{Information refinement preorder}

\begin{definition}[Refinement preorder]
\label{def:preorder}
For two observation maps $O_a, O_b : \Hsym \to S_a, S_b$, write
$O_a \preceq O_b$ (``$O_b$ refines $O_a$'') if there is a function $f$
such that $O_a = f \circ O_b$ on $\Hsym$.  Write $O_a \bowtie O_b$
for \emph{incomparability}: $O_a \npreceq O_b$ and $O_b \npreceq O_a$.
\end{definition}

This is a preorder (reflexive and transitive) but not in general
antisymmetric.  No lattice structure is claimed.  Note that
$O_a \preceq O_b$ is equivalent to: every $O_b$-collision is also an
$O_a$-collision; or equivalently, $O_b$ has finer fibres than $O_a$.

\subsection{Building blocks: the on-line pair and the off-line quartet}

Two basic building blocks recur throughout:
\begin{align*}
  L(t) &= \{{\textstyle\frac{1}{2}} + it,\; {\textstyle\frac{1}{2}} - it\}
         && \text{(critical-line conjugate pair, }t > 0\text{)},\\
  Q(\sigma_0, T) &= \{\sigma_0 \pm iT,\; 1-\sigma_0 \pm iT\}
         && \text{(off-line symmetric quartet, }\sigma_0 \neq {\textstyle\frac{1}{2}}\text{)}.
\end{align*}
Both are in $\Hsym$ for any $0 < \sigma_0 < 1$ and $t, T > 0$.
We have $P(L(t)) = 1$ and $P(Q(\sigma_0,T)) = 0$ (for $\sigma_0 \neq 1/2$).

Direct computation using Definition~\ref{def:Li} gives:
\begin{align}
  \Li_j(L(t)) &= 4\,\re\varphi_j\!\bigl({\textstyle\frac12}+it\bigr),
                \label{eq:Lionline}\\
  \Li_j(Q(\sigma_0,T))
    &= 4\,\re\bigl[\varphi_j(\sigma_0+iT) + \varphi_j(1-\sigma_0+iT)\bigr].
    \label{eq:Lioffl}
\end{align}
For the standard Li tests, direct substitution gives
\[
  \Li_j(Q(\sigma_0,T))
  = \frac{4j^2 \cos\bigl(j\arctan(T/\sigma_0)+j\arctan(T/(1-\sigma_0))\bigr)}
         {T^2\bigl((\sigma_0^2+T^2)((1-\sigma_0)^2+T^2)\bigr)^{j/2}}
    + O(T^{-4}).
\]
In particular, $\Li_j(Q(\sigma_0,T)) \to 0$ as $T\to\infty$.

%%--------------------------------------------------------------------
\section{Li observations admit exact collisions (Theorem~1)}
\label{sec:B2}
%%--------------------------------------------------------------------

\subsection{Statement}

\begin{theorem}[Exact Li collision]
\label{thm:B2}
For every $m \geq 1$ there exist $\mathcal{Z}_+, \mathcal{Z}_- \in \Hsym$
with $P(\mathcal{Z}_+) = 1$, $P(\mathcal{Z}_-) = 0$, and
\[
  \Li_j(\mathcal{Z}_+) = \Li_j(\mathcal{Z}_-) \qquad
  \text{for all } j = 1, \ldots, m,
\]
where $\Li_j$ uses the standard Li tests $\varphi_j(\rho)=1-(1-\rho^{-1})^j$.
Moreover, $\mathcal{Z}_\pm$ can be taken to have rational heights and all
multiplicities nonneg integers.
\end{theorem}

\subsection{The collision matrix}

Fix distinct positive rationals $0 < t_1 < \cdots < t_m$ and a rational
$T > t_m$.  Let $\mathcal{Z}_+$ consist of $M$ copies of each $L(t_k)$
(for a buffer $M \geq 1$ to be chosen), and let $\mathcal{Z}_-$ be obtained
from $\mathcal{Z}_+$ by replacing the multiplicities: $\mathcal{Z}_-$
contains $(M + n_k)$ copies of $L(t_k)$ and $R$ copies of $Q(3/4, T)$,
where $n_k \in \ZZ$ with $M + n_k \geq 0$ are to be determined.

The exact collision condition $\Li_j(\mathcal{Z}_+) = \Li_j(\mathcal{Z}_-)$
becomes, for $j = 1, \ldots, m$:
\[
  \sum_{k=1}^m n_k \, C_{jk} + R \, d_j(T) = 0,
  \tag{$\star$}\label{eq:collision}
\]
where
\begin{align*}
  C_{jk} &= \Li_j(L(t_k)) = 4\,\re\varphi_j\!\bigl({\textstyle\frac12}+it_k\bigr),\\
  d_j(T) &= \Li_j(Q(3/4,T))
           = 4\,\re\bigl[\varphi_j({\textstyle\frac34}+iT)
             + \varphi_j({\textstyle\frac14}+iT)\bigr].
\end{align*}
In matrix form, $Cn + Rd(T) = 0$ where $C \in \RR^{m \times m}$
and $d(T) \in \RR^m$.

\subsection{Non-vanishing of the quartet tail}

\begin{lemma}[Positivity of $d_1(T)$]
\label{lem:d1pos}
For the standard Li test $\varphi_1(\rho)=\rho^{-1}$ and every $T > 0$:
$d_1(T) > 0$.
\end{lemma}

\begin{proof}
Direct computation:
\begin{align*}
  d_1(T) &= 4\,\re\bigl[\tfrac{1}{\frac34+iT} + \tfrac{1}{\frac14+iT}\bigr]
           = 4\biggl(\frac{3/4}{(3/4)^2+T^2} + \frac{1/4}{(1/4)^2+T^2}\biggr)\\
         &= \frac{64(16T^2+3)}{256T^4+160T^2+9} > 0,
\end{align*}
since both numerator and denominator have only positive coefficients.
\end{proof}

\subsection{Non-singularity of the Li collision matrix}

\begin{lemma}[Full rank of $C$ over $\QQ$]
\label{lem:Crank}
For distinct positive rationals $t_1, \ldots, t_m$, the matrix
$C = (C_{jk})$ with $C_{jk} = 4(1 - T_j(x_k))$
\textup{(}where $x_k = (4t_k^2 - 1)/(4t_k^2+1)$ and $T_j$ is the
degree-$j$ Chebyshev polynomial\textup{)} lies in $\QQ^{m\times m}$
and satisfies $\det C \neq 0$.
\end{lemma}

\begin{proof}
Since $t_k \in \QQ$, each $x_k \in \QQ$; since $T_j$ has integer
coefficients, $C_{jk} = 4(1-T_j(x_k)) \in \QQ$.  For non-singularity:
$T_j(1) = 1$ gives $1 - T_j(x) = (1-x)\,q_j(x)$ for a polynomial $q_j$
of degree exactly $j-1$ with leading coefficient $2^{j-1}$.  Thus
\[
  C_{jk} = 4(1-x_k)\,q_j(x_k),
\]
and $C = 4 \cdot [q_j(x_k)] \cdot \operatorname{diag}(1-x_k)$.
Since $x_k \in (-1,1)$, each $1-x_k > 0$, so $\det C \neq 0$ iff
$\det[q_j(x_k)] \neq 0$.  The polynomials $\{q_j\}_{j=1}^m$ have
degrees $0,1,\ldots,m-1$ with positive leading coefficients, so
$[q_j(x_k)] = A \cdot V(x_1, \ldots, x_m)$ where $A$ is lower-triangular
with positive diagonal entries and $V$ is the Vandermonde matrix.
The $x_k$ are distinct (since $x(t)=(4t^2-1)/(4t^2+1)$ is strictly
increasing), so $\det V \neq 0$.  Hence $\det C \neq 0$.
\end{proof}

\begin{remark}
\label{rem:momentVander}
For moment-type tests $\varphi_j(\rho) = \rho^{-j}$, the matrix entries
$C_{jk} = 4\re(1/2+it_k)^{-j}$ satisfy a cosine-Vandermonde relation in
the variable $\cos\theta_k$ with $\theta_k = 2\arctan(2t_k)$, and an
analogous rank argument applies.
\end{remark}

\subsection{Rational solution and nonneg integer multiplicities}

\begin{proof}[Proof of Theorem~\ref{thm:B2}]
Choose distinct rational $t_1 < \cdots < t_m$ and rational $T > t_m$.
By Lemma~\ref{lem:Crank}, $C \in \QQ^{m\times m}$ and $\det C \neq 0$.
When $T \in \QQ$, each $d_j(T)$ is a rational number (sum of rational
functions of $T$ evaluated at rational~$T$), so $d(T) \in \QQ^m$.  Set
\[
  \beta = -C^{-1} d(T) \in \QQ^m.
\]
Let $R = \operatorname{lcm}(\text{denominators of } \beta_1,\ldots,\beta_m)
\in \ZZ_{\geq 1}$ and $n_k = R\beta_k \in \ZZ$.
Then $Cn + Rd(T) = 0$, i.e.\ \eqref{eq:collision} holds.

\textbf{Non-negative multiplicities.} Set $M = \max_k |n_k|$.  Then
$M + n_k \geq 0$ for all~$k$.  Let $\mathcal{Z}_+$ contain $M$ copies
of $L(t_k)$ for each $k = 1,\ldots,m$, and let $\mathcal{Z}_-$ contain
$(M + n_k)$ copies of $L(t_k)$ plus $R$ copies of $Q(3/4, T)$.
Both are in $\Hsym$; all multiplicities are $\geq 0$.

\textbf{Observation collision.} By construction $Cn + Rd(T) = 0$, which is
exactly $\Li_j(\mathcal{Z}_-) - \Li_j(\mathcal{Z}_+) = 0$ for each~$j$.

\textbf{Predicate separation.}  $P(\mathcal{Z}_+) = 1$ (all atoms on
$\re\rho = 1/2$).  $P(\mathcal{Z}_-) = 0$ since $\mathcal{Z}_-$ contains
atoms at $\re\rho = 3/4$.  This requires $R \geq 1$: indeed
$d_1(T) > 0$ (Lemma~\ref{lem:d1pos}) implies $d(T) \neq 0$, hence
$\beta = -C^{-1}d(T) \neq 0$, hence $R \geq 1$.
\end{proof}

\subsection{Numerical illustration ($m=2$, exact rational arithmetic)}

We exhibit the complete collision for $m=2$ with parameters
$t_1=1$, $t_2=2$, $T=3$, $\sigma_0=\frac{3}{4}$.

\textbf{Chebyshev nodes.}
$x_k = (4t_k^2-1)/(4t_k^2+1)$:
\[
  x_1 = \tfrac{3}{5}, \qquad x_2 = \tfrac{15}{17}.
\]

\textbf{Collision matrix.}  $C_{jk} = 4(1-T_j(x_k))$, where
$T_1(x)=x$, $T_2(x)=2x^2-1$:
\[
  C = \begin{pmatrix} \tfrac{8}{5} & \tfrac{8}{17} \\[4pt]
                      \tfrac{128}{25} & \tfrac{512}{289} \end{pmatrix},
  \qquad \det C = \tfrac{3072}{7225} \neq 0.
\]

\textbf{Quartet tail.}  $d_j = \Li_j(Q(3/4, 3))$ via \eqref{eq:Lioffl}:
\[
  d_1 = \tfrac{3136}{7395}, \qquad d_2 = \tfrac{10041088}{6076225}.
\]

\textbf{Rational solution.}  $\beta = -C^{-1}d$:
\[
  \beta_1 = \tfrac{7978}{128673}, \qquad
  \beta_2 = -\tfrac{42082}{37845}.
\]
Clearing denominators: $R = \mathrm{lcm}(128673,\, 37845) = 643365$, giving
\[
  n_1 = R\beta_1 = 39890, \qquad n_2 = R\beta_2 = -715394.
\]
Buffer: $M = 715394$, so $M+n_1 = 755284 \geq 0$ and $M+n_2 = 0 \geq 0$.

\textbf{Exact zero residual} (verified by exact rational arithmetic with
\texttt{fractions.Fraction}):
\[
  C_{j1}n_1 + C_{j2}n_2 + R\,d_j = 0 \quad \text{for } j=1,2.
\]

\textbf{The collision pair.}  $\mathcal{Z}_+$ contains $M = 715394$ copies
of $L(1)$ and $L(2)$; $\mathcal{Z}_-$ contains $755284$ copies of $L(1)$,
$0$ copies of $L(2)$, and $R = 643365$ copies of $Q(3/4, 3)$.
Both are in $\Hsym$ with nonneg integer multiplicities;
$\Li_j(\mathcal{Z}_+) = \Li_j(\mathcal{Z}_-)$ exactly for $j=1,2$;
and $P(\mathcal{Z}_+) = 1$, $P(\mathcal{Z}_-) = 0$.

%%--------------------------------------------------------------------
\section{No uniform Li separation (Theorem~2)}
\label{sec:B1}
%%--------------------------------------------------------------------

\subsection{Statement}

\begin{theorem}[No uniform separation]
\label{thm:B1}
Fix $m \geq 1$.  Let $\mathcal{Z}_+ \in \Hsym$ satisfy $P(\mathcal{Z}_+)=1$
and suppose each $\Li_j(\mathcal{Z}_+)$ strictly exceeds a threshold
$c_j \in \RR$.  Then for every $\varepsilon > 0$ there exists
$\mathcal{Z}_- \in \Hsym$ with $P(\mathcal{Z}_-) = 0$ and
\[
  |\Li_j(\mathcal{Z}_-) - \Li_j(\mathcal{Z}_+)| < \varepsilon \qquad
  \text{for all } j = 1, \ldots, m.
\]
In particular, $\mathcal{Z}_-$ also satisfies $\Li_j(\mathcal{Z}_-) > c_j$
whenever $\varepsilon < \Li_j(\mathcal{Z}_+) - c_j$.
\end{theorem}

\begin{proof}
Add a single off-line quartet $Q(3/4, T)$ to $\mathcal{Z}_+$ to obtain
$\mathcal{Z}_- = \mathcal{Z}_+ \sqcup Q(3/4, T)$.  Then
\[
  \Li_j(\mathcal{Z}_-) = \Li_j(\mathcal{Z}_+) + d_j(T),
\]
where $d_j(T) = \Li_j(Q(3/4,T)) \to 0$ as $T \to \infty$ for the standard
Li tests (since $d_j(T) = O(T^{-2})$ from the quartet formula).
Hence for $T$ large enough, $|\Li_j(\mathcal{Z}_-) - \Li_j(\mathcal{Z}_+)| < \varepsilon$
for all~$j$.  Since $Q(3/4,T)$ contributes atoms at $\re\rho = 3/4 \neq 1/2$,
$P(\mathcal{Z}_-) = 0$.
\end{proof}

\begin{remark}[Approximate vs.\ exact obstruction]
\label{rem:approxvsexact}
Theorem~2 produces an \emph{approximate} non-discrimination: for every
$\varepsilon>0$ one can find $\mathcal{Z}_-$ with observations within
$\varepsilon$ of $\mathcal{Z}_+$, but not necessarily equal.  This is not
an information obstruction in the exact sense of Section~\ref{sec:intro}
(which requires $O(\mathcal{Z}_+)=O(\mathcal{Z}_-)$ exactly).  Theorem~1
provides the exact version; Theorem~2 is complementary in showing that the
gap cannot be uniformly bounded away from zero.
\end{remark}

\begin{corollary}[Positivity threshold]
\label{cor:B1Li}
For every fixed $K$ and every $\mathcal{Z}_+ \in \Hsym$ with $P(\mathcal{Z}_+)=1$
and $\Li_j(\mathcal{Z}_+) > 0$ for $j=1,\ldots,K$, there exists
$\mathcal{Z}_- \in \Hsym$ with $P(\mathcal{Z}_-) = 0$ and
$\Li_j(\mathcal{Z}_-) > 0$ for $j=1,\ldots,K$.
\end{corollary}

\begin{proof}
Any $\mathcal{Z}_+ = L(t_0)$ for rational $t_0 > 0$ satisfies
$P(\mathcal{Z}_+)=1$ and
$\Li_j(L(t_0)) = 4\,\re\varphi_j(1/2+it_0) > 0$ for all $j \geq 1$
(since $\varphi_j(1/2+it_0)$ has positive real part by direct computation).
Apply Theorem~2 with $\varepsilon = \min_j \Li_j(\mathcal{Z}_+)/2$ to
obtain $\mathcal{Z}_-$ with $P(\mathcal{Z}_-)=0$ and
$\Li_j(\mathcal{Z}_-) > \Li_j(\mathcal{Z}_+)/2 > 0$.
\end{proof}

\subsection{Relation to Theorem~1}

Theorem~2 gives approximate non-discrimination (tolerance $\varepsilon$);
Theorem~1 gives exact non-discrimination ($\varepsilon = 0$).  Theorem~2
requires only that $d_j(T) \to 0$, which holds for the standard Li tests.
Theorem~1 additionally uses the $\QQ$-rationality of the collision matrix
(Lemma~\ref{lem:Crank}).  Both escape routes are identical: an infinite
hierarchy $m \to \infty$, an Euler-product constraint, or the full
functional equation.

%%--------------------------------------------------------------------
\section{Li and counting observations are incomparable (Theorem~3)}
\label{sec:H}
%%--------------------------------------------------------------------

\subsection{Statement}

\begin{theorem}[Incomparability]
\label{thm:Hincomp}
Fix $M \geq 1$ and positive sampling levels $0 < d_1 < \cdots < d_M$.
For every $K \geq 1$,
\[
  \Ofinite^{(K)} \bowtie \Otheta \qquad \text{on } \Hsym.
\]
That is, $\Ofinite^{(K)} \npreceq \Otheta$ and $\Otheta \npreceq \Ofinite^{(K)}$.
\end{theorem}

\begin{remark}[Scope]
Theorem~3 is a statement about two observation maps on $\Hsym$ and the
refinement preorder between them.  It does not directly involve the target
predicate $P$: incomparability of two observation maps means each can
distinguish some pair that the other cannot, not that either alone obstructs
any proof strategy for $P$.  Combining both observations gives an observation
map strictly finer than either alone; such a combined observation might or
might not admit an obstruction.
\end{remark}

\subsection{Witness 1: same-height quartets ($\Ofinite^{(K)} \npreceq \Otheta$)}

For $T > 0$ and $\sigma_0 \in (1/2, 1)$, the multiset $Q(\sigma_0, T)$
has imaginary parts $\pm T$ (each with multiplicity~$2$), so
$N_{Q(\sigma_0,T)}(u) = 4$ for $u \geq T$ and $= 0$ for $u < T$.
In particular, $\Otheta(Q(\sigma_0, T))$ depends only on~$T$, not on
$\sigma_0$.

Taking $T = 10$, $\sigma_0 = 3/4$ and $\sigma_0' = 9/10$:
$\Otheta(Q(3/4, 10)) = \Otheta(Q(9/10, 10))$ at every level $d_m$.

However, by Definition~\ref{def:Li}:
\begin{align*}
  \Li_1(Q(3/4,10)) &= 4\,\re\bigl[\tfrac{1}{3/4+10i}
                      + \tfrac{1}{1/4+10i}\bigr]
                    = \frac{102592}{2576009} \approx 0.039826,\\
  \Li_1(Q(9/10,10)) &= 4\,\re\bigl[\tfrac{1}{9/10+10i}
                       + \tfrac{1}{1/10+10i}\bigr]
                     = \frac{4003600}{100820081} \approx 0.039710.
\end{align*}
The exact difference is $\frac{30024117552}{259713436036729} > 0$
(verified by exact rational arithmetic).

Therefore the pair $(Q(3/4,10),\, Q(9/10,10))$ has
$\Otheta(Q(3/4,10)) = \Otheta(Q(9/10,10))$ but
$\Li_1(Q(3/4,10)) \neq \Li_1(Q(9/10,10))$,
proving $\Ofinite^{(K)} \npreceq \Otheta$.

\subsection{The non-cancellation lemma}

\begin{lemma}[Non-cancellation parameter selection]
\label{lem:noncancel}
Fix $K \geq 1$ and $d_* > 0$.  One can choose rational parameters
$0 < t_1 < \cdots < t_K < T < d_*$ such that the Theorem~1 vector
$\beta = -C^{-1}d(T) \in \QQ^K$ satisfies
\[
  G_K(t_1, \ldots, t_K, T) := 2 + \sum_{k=1}^K \beta_k \neq 0.
\]
\end{lemma}

\begin{proof}
For fixed rational $t_1 < \cdots < t_K$, the entries of $C$, $d(T)$, and
$\beta = -C^{-1}d(T)$ are rational functions of $T$ (since all ingredients
are rational functions of the input $T$ evaluated at rational $t_k$).
Thus $G_K = 2 + \sum_k \beta_k$ is a rational function of
$(t_1,\ldots,t_K,T)$ on the open region
$\Omega = \{0 < t_1 < \cdots < t_K < T\}$ (no upper bound on $T$).

As $T \to \infty$ with $t_k$ fixed, $d_j(T) = O(T^{-2}) \to 0$, so
$\beta \to 0$ and $G_K \to 2 \neq 0$.  Hence $G_K$ is not identically
zero on $\Omega$.  By the identity theorem for real-analytic functions,
$G_K$ is therefore not identically zero on any connected open subset of
$\Omega$.  In particular, $G_K \not\equiv 0$ on
$\Omega_{d_*} = \{0 < t_1 < \cdots < t_K < T < d_*\}$, which is a
non-empty connected open subset of $\Omega$.  The set where $G_K \neq 0$
is then a dense open subset of $\Omega_{d_*}$; since rational points are
dense, one can choose all parameters rational with $G_K \neq 0$.
\end{proof}

\begin{remark}
The lemma is necessary: the exact rational counterexample $K=2$,
$t_1 = 1/24$, $t_2 = 9/40$, $T = 3/8$ gives $G_2 = 0$
(verified by exact computation).  The claim that any Theorem~1 pair
automatically separates the counting observation is false; one must
select parameters avoiding the cancellation locus.
\end{remark}

\subsection{Witness 2: Theorem~1 pair with non-cancellation
  ($\Otheta \npreceq \Ofinite^{(K)}$)}

\begin{proof}[Proof of $\Otheta \npreceq \Ofinite^{(K)}$]
Fix $K \geq 1$.  Choose rational $0 < t_1 < \cdots < t_K < T < d_1$ using
Lemma~\ref{lem:noncancel}, so that $G_K \neq 0$.

Construct the Theorem~1 collision pair: let $R = \operatorname{lcm}(\text{denom of }\beta)$,
$n_k = R\beta_k$, $M = \max_k |n_k|$.  Set $\mathcal{Z}_+$ with $M$ copies
of $L(t_k)$ and $\mathcal{Z}_-$ with $(M+n_k)$ copies of $L(t_k)$ plus
$R$ copies of $Q(3/4,T)$.  By Theorem~\ref{thm:B2},
$\Ofinite^{(K)}(\mathcal{Z}_+) = \Ofinite^{(K)}(\mathcal{Z}_-)$.

For the counting observation, since $T < d_1$, all atoms of both
$\mathcal{Z}_+$ and $\mathcal{Z}_-$ have $|\im\rho| \leq T < d_1$.
Computing:
\begin{align*}
  N_{\mathcal{Z}_-}(d_1) - N_{\mathcal{Z}_+}(d_1)
  &= 2\sum_{k=1}^K n_k + 4R
   = 2R\!\left(\sum_{k=1}^K \beta_k + 2\right)
   = 2R \, G_K.
\end{align*}
Since $G_K \neq 0$ and $R \geq 1$, we have
$N_{\mathcal{Z}_-}(d_1) \neq N_{\mathcal{Z}_+}(d_1)$, so
$\Otheta(\mathcal{Z}_-) \neq \Otheta(\mathcal{Z}_+)$, proving
$\Otheta \npreceq \Ofinite^{(K)}$.
\end{proof}

\subsection{Proof of Theorem~\ref{thm:Hincomp}}

Both directions are established:
$\Ofinite^{(K)} \npreceq \Otheta$ by Witness~1 (Section~\ref{sec:H}.2)
and $\Otheta \npreceq \Ofinite^{(K)}$ by Witness~2 (Section~\ref{sec:H}.4).
This gives $\Ofinite^{(K)} \bowtie \Otheta$ for every $K \geq 1$. \qed

\subsection{Remark on the oracle}

The finest observation is $O_{\mathrm{oracle}}(\mathcal{Z}) = \mathcal{Z}$
(the full multiset).  Both $\Ofinite^{(K)}$ and $\Otheta$ are strictly
coarser: the Theorem~1 pair exhibits
$\Ofinite^{(K)}(\mathcal{Z}_+) = \Ofinite^{(K)}(\mathcal{Z}_-)$ with
$\mathcal{Z}_+ \neq \mathcal{Z}_-$, and similarly for $\Otheta$ via the
same-height quartet.

%%--------------------------------------------------------------------
\section{Finite Euler factors do not force critical-line zeros (Corollary~4)}
\label{sec:C}
%%--------------------------------------------------------------------

\subsection{Setting: Helson zeta functions}

For a completely multiplicative unimodular function
$\chi : \mathbb{N} \to \TT$ ($|\chi(p)| = 1$ for all primes $p$), define
the Helson zeta function
\[
  \zeta_\chi(s) = \sum_{n \geq 1} \chi(n) n^{-s}
               = \prod_p (1-\chi(p)p^{-s})^{-1},
\]
convergent for $\re s > 1$.  For a finite prime cutoff $P_0$, say $\chi$
is \emph{$P_0$-standard} if $\chi(p) = 1$ for all primes $p \leq P_0$.

Let $\mathcal{H}_S$ denote the Helson zeta functions admitting meromorphic
continuation to the open strip $S = \{0 < \re s < 1\}$.  Define the
target predicate $P_S(\zeta_\chi) = 1$ iff every zero $\rho \in S$ of
$\zeta_\chi$ satisfies $\re\rho = 1/2$.

\subsection{Statement}

\begin{corollary}[Finite Euler factors do not force critical-line zeros]
\label{cor:C}
For every finite prime cutoff $P_0$ and every $z_1 \in S$ with
$\re z_1 \neq 1/2$, there exist:
\begin{enumerate}[label=\textup{(\alph*)}]
  \item a completely multiplicative unimodular $\tilde\chi$ that is
    $P_0$-standard, such that $\zeta_{\tilde\chi} \in \mathcal{H}_S$
    and $P_S(\zeta_{\tilde\chi}) = 0$;
  \item a completely multiplicative unimodular $\chi^+$ that is also
    $P_0$-standard, such that $\zeta_{\chi^+} \in \mathcal{H}_S$
    and $P_S(\zeta_{\chi^+}) = 1$.
\end{enumerate}
In particular, the pair $(\zeta_{\chi^+}, \zeta_{\tilde\chi})$ is an
information obstruction pair for the finite-Euler-factor observation
$O_{P_0}(\zeta_\chi) = (\chi(p))_{p \leq P_0}$: equal observations,
opposite values of $P_S$.
\end{corollary}

\begin{proof}
\textbf{Part (a): constructing $P_S=0$.}

\textit{Step 1.}  By Andersson~\cite[Theorem~5]{Andersson2024}, applied
with $U = \CC$, $Z = \{z_1\}$ (a single prescribed simple zero),
there exists a completely multiplicative unimodular $\chi$ such that
$\zeta_\chi$ extends meromorphically to $S$ with a simple zero at $z_1$
(and $\re z_1 \neq 1/2$, so $P_S(\zeta_\chi)=0$).

\textit{Step 2.}  Define
\[
  \tilde\chi(p) =
    \begin{cases} 1 & p \leq P_0, \\ \chi(p) & p > P_0, \end{cases}
\]
extended completely multiplicatively to $\mathbb{N}$.  Then $\tilde\chi$ is
completely multiplicative, unimodular, and $P_0$-standard.
For $\re s > 1$, the identity
\[
  \zeta_{\tilde\chi}(s) = \zeta_\chi(s) \cdot
    \underbrace{\prod_{p \leq P_0} \frac{1 - \chi(p)p^{-s}}{1 - p^{-s}}}_{=:R(s)}
\]
holds by rearranging Euler factors.  Since $|\chi(p)| = |p^{-s}| = p^{-\re s}$
and each factor of $R$ is a ratio of first-degree polynomials in $p^{-s}$
with zeros on the line $\re s = 0$, $R$ extends to a function holomorphic
and nonvanishing on all of $S$.  Continuing both sides meromorphically to
$S$ gives $\zeta_{\tilde\chi} = \zeta_\chi \cdot R$ on $S$.
Since $R(z_1) \neq 0$, the zero at $z_1$ is preserved, so
$P_S(\zeta_{\tilde\chi}) = 0$.

\textbf{Part (b): constructing $P_S=1$.}
The trivial character $\chi^+ \equiv 1$ gives $\zeta_{\chi^+} = \zeta$,
the Riemann zeta function; this is $P_0$-standard and has all nontrivial
zeros on the critical line if and only if RH holds.  For an unconditional
example: take $\chi^+ \equiv 1$ restricted to primes $p > P_0$ (extended
completely multiplicatively with $\chi^+(p)=1$ for all $p$, i.e.\
$\chi^+ \equiv 1$), so $\zeta_{\chi^+} = \zeta$.  Alternatively, apply
Andersson~\cite[Theorem~5]{Andersson2024} with a prescribed zero
$z_1' = 1/2 + i\gamma$ (on the critical line), yielding a Helson $\zeta$
with $P_0$-standard Euler factors and $P_S = 1$.  This makes
$(\zeta_{\chi^+}, \zeta_{\tilde\chi})$ an information obstruction pair
for $O_{P_0}$.
\end{proof}

\begin{remark}[Scope limitation]
Corollary~4 applies only to the Helson class.  It does \emph{not} apply
to the Selberg class, which requires the functional equation, the gamma
factor, a finite analytic conductor, and Ramanujan-type bounds on
coefficients --- none of which are satisfied by a generic Helson zeta
function.  In particular, the multiplicativity of $\chi$ does not make
$\zeta_\chi$ a Selberg-class element.  The escape route for Corollary~4
is any proof method that uses the full Selberg-class structure or the
specific analytic properties of $\zeta$.
\end{remark}

%%--------------------------------------------------------------------
\section{Discussion}
\label{sec:discussion}
%%--------------------------------------------------------------------

\subsection{What the results say together}

The three main theorems cover complementary aspects of the
finite-arithmetic landscape within $\Hsym$:

\begin{itemize}
\item Theorems~1 and~2 together say: \emph{no fixed finite collection of
  standard Li values, however large, can determine $P$ on $\Hsym$.}
  Theorem~1 achieves exact non-determination; Theorem~2 says the gap
  cannot be bounded away from zero.

\item Theorem~3 says: \emph{Li values and counting data are not
  redundant---each can distinguish inputs the other cannot.}  The
  incomparability is a property of the two observation maps alone, not
  a direct obstruction for $P$.  Two observations that are individually
  non-determining may or may not determine $P$ jointly; that question is
  open (see Open problem~2 below).

\item Corollary~4 says: \emph{finite Euler-product information leaves a gap
  within the Helson class}, and exhibits an explicit information obstruction
  pair.
\end{itemize}

A complete finite-arithmetic proof of RH targeting $\Hsym$ would need to
use all three types of information simultaneously.

\subsection{Escape routes}

The following methods are \emph{not} blocked by the results in this paper:

\begin{enumerate}
\item An unbounded hierarchy of Li tests ($m \to \infty$, all simultaneously).
\item The full Euler product $\prod_p (1-p^{-s})^{-1}$ plus functional equation.
\item Selberg-class axioms (for Corollary~4 only).
\item Methods that use the imaginary part of $\log\zeta(1/2+it)$ continuously.
\item The combined observation $\Ofinite^{(K)} \times \Otheta$ (incomparability
  of the two maps does not rule out joint determination of $P$).
\end{enumerate}

\subsection{Open problems}

\begin{enumerate}
\item \textbf{Grow-with-$T$ obstruction.} Is there an information
  obstruction that remains valid even as $K = K(T)$ grows with the height
  parameter $T$?

\item \textbf{Combined Li + counting obstruction.} Does combining
  $\Ofinite^{(K)}$ with $\Otheta$ still leave an information gap for $P$
  on $\Hsym$, or does the combined observation determine $P$ for
  sufficiently large $K, M$?  (The ambient class $\Hsym$ serves as a
  common domain for both maps.)

\item \textbf{Weil-type observations.}
  For a fixed $h \in C_c^\infty(\RR)$, the Weil-type observation
  $\varphi_h(\rho) = \hat{h}((\rho-\tfrac12)/i)$ is not meromorphic and
  falls outside the scope of Theorem~1.  Does an analogue of
  Theorem~1 hold for such observations?  The Chebyshev--Vandermonde
  reduction does not directly apply; a separate argument is needed.

\item \textbf{Information hierarchy convergence.}  Is there a natural
  sequence of observations $O_1 \preceq O_2 \preceq \cdots$ converging to
  $O_{\mathrm{oracle}}$ in some appropriate sense, and at what level does
  the chain separate $P = 1$ from $P = 0$?
\end{enumerate}

%%--------------------------------------------------------------------
\bibliographystyle{amsalpha}
\begin{thebibliography}{99}

\bibitem{Andersson2024}
J.~Andersson,
\emph{Mittag-Leffler type theorems for Helson zeta-functions},
arXiv:2408.15713v1, 2024.

\bibitem{BakerGillSolovay1975}
T.~Baker, J.~Gill, and R.~Solovay,
\emph{Relativizations of the $P=?NP$ question},
SIAM J. Comput.~\textbf{4} (1975), no.~4, 431--442.

\bibitem{BochkovRomanov2022}
I.~Bochkov and R.~Romanov,
\emph{On zeroes and poles of Helson zeta functions},
J.\ Funct.\ Anal.~\textbf{282} (2022), no.~1, 109398.

\bibitem{Bombieri2000}
E.~Bombieri,
\emph{Remarks on Weil's quadratic functional in the theory of prime
numbers I},
Rend.\ Mat.\ Accad.\ Lincei (9)~\textbf{11} (2000), 183--233.

\bibitem{BombieriLagarias1999}
E.~Bombieri and J.~C.~Lagarias,
\emph{Complements to Li's criterion for the Riemann hypothesis},
J.\ Number Theory~\textbf{77} (1999), 274--287.

\bibitem{CCM2025}
A.~Connes, C.~Consani, and H.~Moscovici,
\emph{Zeta spectral triples},
arXiv:2511.22755, 2025.

\bibitem{ConnesWeilFF1999}
A.~Connes,
\emph{Trace formula in noncommutative geometry and the zeros of the
Riemann zeta function},
Selecta Math.\ (N.S.)~\textbf{5} (1999), no.~1, 29--106.

\bibitem{Helson1969}
H.~Helson,
\emph{Compact groups and Dirichlet series},
Ark.\ Mat.~\textbf{8} (1969), 139--143.

\bibitem{Hejhal1976}
D.~A.~Hejhal,
\emph{The Selberg Trace Formula for $\mathrm{PSL}(2,\mathbb{R})$},
Vol.~1, Lecture Notes in Math.~548, Springer, Berlin, 1976.

\bibitem{Li1997}
X.-J.~Li,
\emph{The positivity of a sequence of numbers and the Riemann hypothesis},
J.\ Number Theory~\textbf{65} (1997), 325--333.

\bibitem{RazborovRudich1997}
A.~Razborov and S.~Rudich,
\emph{Natural proofs},
J.\ Comput.\ System Sci.~\textbf{55} (1997), no.~1, 24--35.

\bibitem{Seip2020}
K.~Seip,
\emph{Universality and distribution of zeros and poles of some zeta
functions},
J.\ Anal.\ Math.~\textbf{141} (2020), 331--381.

\end{thebibliography}

\end{document}
